\documentclass[11pt]{article}

\usepackage[margin=1in]{geometry}
\usepackage{graphicx}
\usepackage{amsmath}
\usepackage{amssymb}
\usepackage{hyperref}
\usepackage{caption}
\usepackage{subcaption}
\usepackage{float}
\usepackage{natbib}

\hypersetup{
  colorlinks=true,
  linkcolor=blue,
  urlcolor=blue,
  citecolor=blue
}

% ---------- Title ----------
\title{CTA200: Mini Coding Project}
\author{Isabella Rivera \\
  \small SURP - Supervisors: Mattias Lazda, Prof.\ Maria Drout, Prof.\ Juan Mena-Parra}
\date{May 17, 2026}

\begin{document}
\maketitle

\section{Introduction}
\label{sec:intro}

Radio astronomy probes astrophysical emission over a broad range of
frequencies, from roughly 3~kHz to 900~GHz \citep{nrao_radio}. Depending on the frequency at which a radio source is observed, the source
can appear different. A single radio image therefore captures a source at
only one frequency, while multi-frequency observations can reveal
complementary information about its structure that cannot be obtained from
any single observation. In this mini-project, I examine FITS images of the
same source obtained from three radio surveys spanning roughly an order of
magnitude in observing frequency: VLASS
\citep[Very Large Array Sky Survey, 3~GHz;][]{lacy_2020}, FIRST
\citep[Faint Images of the Radio Sky at Twenty centimeters, 1.4~GHz;][]{becker_1995},
and LoTSS \citep[LOFAR Two-metre Sky Survey, 300~MHz;][]{shimwell_2017}.
This range allows the source morphology to be compared across very
different angular resolutions and wavelength regimes.

Using the \texttt{radioquery} package developed by Mattias Lazda, I downloaded the survey images of the source and produced right ascension-declination maps centered on
the target source for each survey. I then compared the apparent morphology of
the source across the three observing frequencies. One important factor
affecting these images is angular resolution. For a radio telescope or interferometer, the characteristic angular
resolution is set by the ratio of the observing wavelength $\lambda$
to the telescope diameter or maximum interferometric baseline $D$,
%
\begin{equation}
  \theta \sim \frac{\lambda}{D},
  \label{eq:resolution}
\end{equation}
%
where smaller values of $\theta$ correspond to finer angular resolution.
For a fixed telescope size or baseline, longer-wavelength,
lower-frequency observations produce a larger synthesized beam and
therefore a coarser image \citep{condon_ransom_2016}.

In this report, I discuss the methods used in this data analysis in Section~\ref{sec:methods}, present the resulting plots in Section~\ref{sec:results}, and interpret these results and summarize my conclusions in Section~\ref{sec:discussion}.

\section{Methods}
\label{sec:methods}

The project is split across two Python scripts. The first,
\texttt{download\_radio\_image.py}, defines a reusable function that
downloads a FITS file for a given (RA, Dec) and survey name using the
\texttt{radioquery} package made by Mattias Lazda. The second script
\texttt{FITS\_images.ipynb} loads those FITS files and produces RA and
Dec plots.

\subsection{Environment setup}

To download the FITS files for this analysis, I generated an SSH key and added it to my GitHub account to enable authenticated access to my project repository. I then installed the \texttt{radioquery} package inside a dedicated \texttt{conda} environment using the dependencies listed in the \texttt{radioquery} README. This environment included the packages needed to query radio survey images, manipulate FITS files, and produce the final plots, including \texttt{astropy},
\texttt{numpy}, and \texttt{matplotlib}. Using a separate \texttt{conda} environment isolates the project dependencies from the system Python installation, which is useful because the larger-scale version of this analysis will eventually be run on a remote machine with more storage and processing capacity than my laptop.

\subsection{Downloading the data}

I then wrote the function \texttt{download\_radio\_image(ra, dec, survey)} to
download a FITS image for a specified sky position and radio survey. The
function takes the target right ascension, declination, and survey name as
inputs. It converts the coordinates into an
\texttt{astropy.coordinates.SkyCoord} object, checks that the requested
survey is valid, and selects the corresponding \texttt{radioquery} query
class: \texttt{VlassQuery} for VLASS, \texttt{FirstQuery} for FIRST, or
\texttt{LotssQuery} for LoTSS. The resulting FITS file is then saved to an output directory specific to that survey.

This wrapper lets all three surveys be called with the same syntax and can be
applied to many sources by looping over RA/Dec pairs. I tested that it works properly using the
example coordinates from the \texttt{radioquery} README, RA~=~10h50m07.270s
and Dec~=~$+30^\circ40'37.52''$.

\subsection{Loading and locating the source}

The plotting script I made, \texttt{FITS\_images.ipynb}, opens each downloaded FITS
file using 
\texttt{astropy.io.fits}. Radio FITS images often store their data with additional dimensions for
frequency and polarization, even when only one frequency channel and one
polarization image are present. These extra single-element dimensions were
removed using \texttt{.squeeze()}, producing a two-dimensional image suitable
for plotting. The celestial world coordinate system is then extracted
from the FITS header using \texttt{WCS(header).celestial}.

The target position is defined as an ICRS \texttt{SkyCoord} object and
converted separately into pixel coordinates for each survey using
\texttt{wcs.world\_to\_pixel}. This step is necessary because each survey has
a different pixel scale and WCS, so the same sky coordinate does not correspond
to the same pixel location in each FITS image.

\subsection{Matching the cutout region across surveys}
\label{ssec:crop}

To compare the source morphology across surveys, each image was cropped to the
same angular size on the sky rather than the same number of pixels. This is
necessary because VLASS, FIRST, and LoTSS have different pixel scales, so the
same pixel width would not correspond to the same sky area in each image.

I therefore used a helper function, \texttt{arcsec\_to\_pixels(wcs, arcsec)},
to convert my chosen angular half-width of $90''$ into the correct number of
pixels for each survey. This conversion uses the WCS pixel scale, since each
survey has a different number of arcseconds per pixel. A half-width of $90''$
on each side of the source gives a total cutout size of $3' \times 3'$,
centered on the target. This ensures that the VLASS, FIRST, and LoTSS panels
all show the same region of sky, so the source morphology can be compared
directly across observing frequencies.

\subsection{Noise and peak estimation}

For each image, the RMS noise was estimated from a source-free region in the top-left corner of the full FITS image (i.e. before the cropping described in Section~\ref{ssec:crop}), selected to avoid the target source and any obvious imaging artifacts. For an example of the noise region selection, see Figure~\ref{fig:noise-region} in Appendix~\ref{app:noise}.

Because the full survey images have different pixel dimensions, the RMS noise was measured using a source-free corner region whose size was chosen separately for each survey. The VLASS image contains a larger number of pixels, allowing a $200 \times 200$ pixel region to be used while still avoiding the target source and obvious artifacts. For the smaller FIRST and LoTSS images, a $30 \times 30$ pixel region was used instead. The peak flux was measured as the maximum pixel value inside the same $3' \times 3'$ cutout region used for the displayed images (as described in Section~\ref{ssec:crop}).

The RMS noise and source peak flux were used to define the colour scale for
each zoomed-in cutout. The lower limit was set to
\texttt{vmin}~=~$-\mathrm{RMS}$ to keep the background noise visible but
subdued, while the upper limit was set to \texttt{vmax} equal to the measured
source peak flux so that the source structure used most of the available
colour range.

\subsection{Plotting choices}

The final plots were displayed using \texttt{imshow}, and the interpolation was set to
\texttt{interpolation='gaussian'} to smooth the visual appearance of the
finite pixel sampling. This does not change the underlying FITS data, but it
makes the plotted images easier to compare visually by reducing the blocky
appearance of individual pixels. Each panel was plotted in right ascension and declination using the celestial
WCS from the corresponding FITS header.

\subsection{Contour maps}

The contour maps were also produced for the
same cutout regions. The contours were plotted using \texttt{ax.contour}. For each survey, the contour levels were set at
$3$, $5$, $10$, $20$, and $50$ times the local RMS noise.

Using contours in units of the RMS noise allows the source structure to be
compared across surveys even though the absolute flux densities and noise
levels differ. The lowest contour, $3\sigma$, traces the faintest significant
emission, while the higher contours show how the emission becomes concentrated
toward the brightest regions of the source.

\section{Results}
\label{sec:results}

Figure~\ref{fig:six-panel} shows the same source in the three different radio surveys,
with both colour-scaled cutouts (left column) and contour maps (right
column).

\begin{figure}[H]
  \centering
  % Row 1: VLASS
  \begin{subfigure}[t]{0.48\textwidth}
    \includegraphics[width=\linewidth]{vlass_panel.pdf}
    \caption{VLASS, 3~GHz - cutout.}
  \end{subfigure}
  \hfill
  \begin{subfigure}[t]{0.48\textwidth}
    \includegraphics[width=\linewidth]{vlass_contour.pdf}
    \caption{VLASS, 3~GHz - contours.}
  \end{subfigure}

  \vspace{0.5em}

  % Row 2: FIRST
  \begin{subfigure}[t]{0.48\textwidth}
    \includegraphics[width=\linewidth]{first_panel.pdf}
    \caption{FIRST, 1.4~GHz - cutout.}
  \end{subfigure}
  \hfill
  \begin{subfigure}[t]{0.48\textwidth}
    \includegraphics[width=\linewidth]{first_contour.pdf}
    \caption{FIRST, 1.4~GHz - contours.}
  \end{subfigure}

  \vspace{0.5em}

  % Row 3: LoTSS
  \begin{subfigure}[t]{0.48\textwidth}
    \includegraphics[width=\linewidth]{lotss_panel.pdf}
    \caption{LoTSS, 300~MHz - cutout.}
  \end{subfigure}
  \hfill
  \begin{subfigure}[t]{0.48\textwidth}
    \includegraphics[width=\linewidth]{lotss_contour.pdf}
    \caption{LoTSS, 300~MHz - contours.}
  \end{subfigure}

  \caption{Cutouts (left column) and contour maps (right column) of the
           same source centered on RA~=~10h50m07.270s,
           Dec~=~$+30^\circ40'37.52''$ in three radio surveys spanning
           roughly a factor of 20 in frequency. Each row shows one
           survey: VLASS at 3~GHz (panels a, b), FIRST at 1.4~GHz
           (panels c, d), and LoTSS at 300~MHz (panels e, f). Contours
           are drawn at $3$, $5$, $10$, $20$, and $50\times$ the local
           RMS noise. Note that the LoTSS peak is well above the
           highest plotted contour ($\sim 280\sigma$), reflecting the
           source's much brighter low-frequency flux density.
           The colour scale in each panel runs from \texttt{vmin}~$=-$RMS to \texttt{vmax}~$=$~peak flux of the source.}
  \label{fig:six-panel}
\end{figure}


\section{Discussion}
\label{sec:discussion}
The source is detected in all three radio surveys, but its apparent morphology changes noticeably with observing frequency. The colour scale images show the overall source structure, while the contour maps provide a clearer comparison of low-level emission by tracing the source relative to the RMS noise in each image. This makes faint emission easier to identify than in the colour-scale images alone.

To start, the images are clearly showing the same object. The overall two component morphology is recognizable in all three surveys as a bright compact source on the bottom left and a nearby, more diffuse ring-like feature to the top right, with the same relative positions and orientation on the sky. This consistency across roughly an order of magnitude in observing frequency is what makes it meaningful to compare the panels in the first place.

The clearest difference between the surveys is angular resolution. The VLASS image shows the sharpest and most compact structure, while the LoTSS image appears broader and more spatially smoothed. This behaviour is consistent with the resolution scaling in Eq.~\ref{eq:resolution}:
for a fixed telescope size or interferometric baseline, longer-wavelength,
lower-frequency observations produce a larger synthesized beam and therefore
a coarser image \citep{condon_ransom_2016}. The same trend appears in the contour maps. The VLASS and FIRST contours show more compact substructure, while the LoTSS contours are smoother and more extended.

The LoTSS image also shows more diffuse emission between the two bright source components. This is most likely because the larger LoTSS beam smooths over small-scale structure, making the emission appear more continuous. In contrast, the higher-frequency VLASS image better resolves compact features, so the source appears sharper but less extended.

The source also appears brightest in the lowest-frequency image. The peak flux density is approximately $0.007$ Jy/beam in VLASS, $0.04$ Jy/beam in FIRST, and $0.14$ Jy/beam in LoTSS, making the LoTSS peak roughly $20$ times larger than the VLASS peak. This indicates that the measured flux density changes with observing frequency.

Finally, the VLASS image shows a grainier background and more low-level $3\sigma$ contours from the noise than FIRST or LoTSS. This likely reflects the sharper angular resolution and smaller synthesized beam of VLASS, which make compact noise features and imaging artifacts more visually apparent. This also explains the vertical stripe-like features visible across the image, particularly in the VLASS contour map. Overall, the comparison demonstrates why multi-frequency radio imaging is useful: each survey emphasizes a different aspect of the same source, from compact high-resolution structure at high frequency to more extended diffuse emission at low frequency.

\appendix
\section{Noise Region Selection}
\label{app:noise}

\begin{figure}[H]
  \centering
  \includegraphics[width=0.8\textwidth]{lotss_noise_diagnostic.pdf}
  \caption{Example of the noise region selection for LoTSS. The cyan
           box marks the source-free corner used to estimate the RMS
           noise, and the green cross marks the source position.}
  \label{fig:noise-region}
\end{figure}

\begin{thebibliography}{9}

\bibitem[Becker et al.(1995)]{becker_1995}
Becker, R.~H., White, R.~L. \& Helfand, D.~J. (1995).
The FIRST Survey: Faint Images of the Radio Sky at Twenty Centimeters.
\emph{The Astrophysical Journal}, 450, 559.
\url{https://ui.adsabs.harvard.edu/abs/1995ApJ...450..559B}.

\bibitem[Condon \& Ransom(2016)]{condon_ransom_2016}
Condon, J.~J. \& Ransom, S.~M. (2016).
\emph{Essential Radio Astronomy}.
Princeton University Press.
Available at
\url{https://www.cv.nrao.edu/~sransom/web/xxx.html}.

\bibitem[Lacy et al.(2020)]{lacy_2020}
Lacy, M. et al. (2020).
The Karl G.\ Jansky Very Large Array Sky Survey (VLASS): Science Case
and Survey Design.
\emph{Publications of the Astronomical Society of the Pacific}, 132, 035001.
\url{https://doi.org/10.1088/1538-3873/ab63eb}.

\bibitem[NRAO()]{nrao_radio}
National Radio Astronomy Observatory.
\emph{The Science of Radio Astronomy}.
\url{https://public.nrao.edu/radio-astronomy/the-science-of-radio-astronomy/}.

\bibitem[Shimwell et al.(2017)]{shimwell_2017}
Shimwell, T.~W. et al. (2017).
The LOFAR Two-metre Sky Survey -- I.\ Survey Description and Preliminary
Data Release.
\emph{Astronomy \& Astrophysics}, 598, A104.
\url{https://doi.org/10.1051/0004-6361/201629313}.

\end{thebibliography}

\end{document}